\documentclass[12pt]{book}
\usepackage[utf8]{inputenc}
\usepackage[T1]{fontenc}
\usepackage{amsmath,amssymb}
\usepackage{calligra}
% Calligraphic font for categories & sheaves (per request)
\newcommand{\Cat}[1]{\mathcal{#1}}
\newcommand{\Sheaf}[1]{\mathcal{#1}}

\begin{document}

\frontmatter

\preface
In the spring of 1963 I suggested to Grothendieck the possibility of my running a seminar at Harvard on his theory of duality for coherent sheaves---a theory which had been hinted at in his talk to the Séminaire Bourbaki in 1958 [8],
and in his talk to the International Congress of Mathematicians in 1958 [9], but had never been developed systematically. He agreed, saying that he would provide an outline of the material, if I would fill in the details and write up lecture notes of the seminar. During the summer of 1963, he wrote a series of ``pré-notes'' [10] which were to be the basis for the seminar.

I quote from the preface of the pré-notes:

``Les présentes notes donnent une esquisse assez détaillée d'une théorie cohomologique de la dualité des Modules cohérents sur les préschémas. Les idées principales de la théorie m'étaient connues dès 1959, mais le manque de fondements adéquats d'Algèbre Homologique m'avait empêché d'aborder une rédaction d'ensemble. Cette lacune de fondements est sur le point d'être comblée par la thèse de VERDIER, ce qui rend en principe possible un exposé satisfaisant. Il est d'ailleurs
apparu depuis qu'il existe des théories cohomologiques de dualité formellement très analogues à celle developpée ici dans toutes sortes d'autres contextes: faisceaux cohérents
sur les espaces analytiques, faisceaux abeliens sur les
espaces topologiques (VERDIER), modules galoisiens (VERDIER, TATE), faisceaux de torsion sur les schémas munis de leur topologie étale, corps de classe en tous genres \ldots Cela me semble une raison assez sérieuse pour se familiariser avec le yoga général de la dualité dans un cas type, comme
la théorie cohomologique des résidus.

La théorie consiste pour l'essentiel dans des questions de variance: construction d'un foncteur $f^!$ et d'un homomorphisme-trace $\mathbf{R}f_*f^! \to \mathrm{id}$. La construction donnée ici est compliquée et indirecte et n'est pas valable sous des conditions aussi générales qu'on est en droit de s'y attendre. Il faudra sans doute une idée nouvelle pour apporter des simplifications substantielles.''

The seminar took place in the fall and winter of 1963--64, with the assistance of David Mumford, John Tate, Stephen Lichtenbaum, John Fogarty, and others, and gave rise to a
series of six exposés which were circulated to a limited audience under the title ``Séminaire Hartshorne''. The present
notes are a revised, expanded and completed version of the previous notes.

I would like to take this opportunity to thank all those people who have helped in the course of this work, and in particular A. Grothendieck, who gave continual support and encouragement throughout the whole project.

\vspace{\baselineskip}
\raggedleft
R.H.\\
Cambridge, May 1966

\endpreface

% Next step: Introduction section ready to add in same style
% All categories (e.g., triangulated categories, derived categories)
% and sheaves will be wrapped in \Cat{} / \Sheaf{} calligraphic font

\end{document}